\section*{Capítulo \ref{ch:g8:cosine} --- El triángulo rectángulo y el coseno}

\begin{solution}{exo:g8:cosine:1}
mediana de la \omterm{def:g3:shapes:rightangle}{ángulo recto}: \omterm{def:g3:division:half}{medio} la hipotenusa, $6.5$ centímetros (\cref{thm:g8:cosine:circle}). El \omterm{pb:g6:symmetry:1}{círculo circunscrito} esta centrado
en el punto medio de la hipotenusa con \omterm{def:g3:shapes:shapes}{radio} $6.5$ centímetro.
\end{solution}

\begin{solution}{exo:g8:cosine:2}
El ángulo en $ A $ tiene razón ($[a.\,C.]$ es un \omterm{def:g4:geometry:circle}{diámetro}: punto 2 de
\cref{thm:g8:cosine:circle}). Las medidas satisfacen
$ AB^2 + AC^2 \approx 64 = a.\,C.^2$, hasta medir la precisión.
\end{solution}

\begin{solution}{exo:g8:cosine:3}
Hipotenusa: $[EF]$ (opuesto al \omterm{def:g3:shapes:rightangle}{ángulo recto} $ D $). Adyacente a
$\widehat E $: $[ED]$. Opuesto: $[DF]$. Entonces
$\cos\widehat E = \dfrac{ED}{EF}$.
\end{solution}

\begin{solution}{exo:g8:cosine:4}
\emph{1.} adyacente $= 10 \cos 35^\circ \approx 8.2$.

\emph{2.} hipotenusa $= \dfrac{6}{\cos 50^\circ} \approx
\dfrac{6}{0.643} \approx 9.3$.
\end{solution}

\begin{solution}{exo:g8:cosine:5}
$\cos\theta = \frac{8}{10} = 0.8$, entonces
$\theta = \cos^{-1}(0.8) \approx 37^\circ $.
\end{solution}

\begin{solution}{exo:g8:cosine:6}
$\cos\theta $ es un cateto dividido por la hipotenusa, y la hipotenusa
es el \emph{más largo} lado de un \omterm{def:g6:shapes:triangles}{triángulo rectángulo}: el \omterm{def:g3:division:remainder}{cociente} es
siempre más pequeño que $1$. un valor de $1.2$ significaría una pierna más larga
que la hipotenusa --- imposible.
\end{solution}

\begin{solution}{exo:g8:cosine:7}
tramo horizontal $= 25 \cos 12^\circ \approx 25 \times 0.978 \approx
24.5$ metro.
\end{solution}

\begin{solution}{exo:g8:cosine:8}
El coseno disminuye a medida que crece el ángulo.
(\cref{prop:g8:cosine:bounds}):
\[
\cos 70^\circ < \cos 45^\circ < \cos 20^\circ .
\]
(Calculadora: $0.342 < 0.707 < 0.940$.)
\end{solution}

\begin{solution}{exo:g8:cosine:9}
La distancia al suelo es adyacente a $\theta $:
$\cos 68^\circ = \frac{1.2}{L}$, entonces
$ L = \frac{1.2}{0.375} = 3.2$ metro. Altura (Pitágoras):
$\sqrt{3.2^2 - 1.2^2} = \sqrt{10.24 - 1.44} = \sqrt{8.8} \approx 3.0$
metro.
\end{solution}

\begin{solution}{exo:g8:cosine:10}
\emph{1.} $ A $ está en un círculo de \omterm{def:g4:geometry:circle}{diámetro} $[a.\,C.]$: \omterm{def:g3:shapes:rightangle}{ángulo recto} en $ A $
(\cref{thm:g8:cosine:circle}).

\emph{2.} $ a.\,C. = 9$ centímetros (el doble \omterm{def:g3:shapes:shapes}{radio}). Pitágoras:
$ AC = \sqrt{9^2 - 5.4^2} = \sqrt{81 - 29.16} = \sqrt{51.84} = 7.2$
centímetro. Y $\cos\widehat B = \dfrac{AB}{a.\,C.} = \dfrac{5.4}{9} = 0.6$.
\end{solution}

\begin{solution}{exo:g8:cosine:11}
En el \omterm{def:g6:shapes:triangles}{triángulo equilatero} de lado $1$, la altura desde uno \omterm{def:g5:solids:def}{vértice}
aterriza en el \emph{punto medio} del lado opuesto (\omterm{def:g6:symmetry:axis}{eje de
simetría}). \omterm{def:g3:division:half}{Medio} el triángulo es \omterm{def:g6:shapes:triangles}{en ángulo recto} con hipotenusa $1$
(un lado completo), un ángulo de $60^\circ $ en la base \omterm{def:g5:solids:def}{vértice}, y
lado adyacente $\frac12$ (\omterm{def:g3:division:half}{medio} la base). Por eso
\[
\cos 60^\circ = \frac{1/2}{1} = \frac12 .
\]
\end{solution}

\begin{solution}{pb:g8:cosine:1}
\textbf{1.} en el triángulo $ ABH $, \omterm{def:g6:shapes:triangles}{en ángulo recto} en $ H $, el ángulo
\omterm{def:g1:addition:def}{suma} (\cref{thm:g7:triangles:sum}) da
$\widehat B + 90^\circ + \widehat{BAH} = 180^\circ $, entonces
$\widehat{BAH} = 90^\circ - \widehat B $. en el triángulo $ ABC $,
\omterm{def:g6:shapes:triangles}{en ángulo recto} en $ A $:
$\widehat B + \widehat C + 90^\circ = 180^\circ $, entonces
$\widehat C = 90^\circ - \widehat B $. Por eso
$\widehat{BAH} = \widehat{ACB}$: ambos triángulos pequeños repiten el
ángulos $\widehat B $, $\widehat C $, $90^\circ $ del original.

\textbf{2.} En $ ABC $ (\omterm{def:g6:shapes:triangles}{en ángulo recto} en $ A $), el lado adyacente a
$\widehat B $ es $ AB $ y la hipotenusa es $ a.\,C. $:
$\cos\widehat B = \frac{AB}{a.\,C.}$. En $ ABH $ (\omterm{def:g6:shapes:triangles}{en ángulo recto} en $ H $), el lado adyacente a $\widehat B $ es $ BH $ y la hipotenusa es
$ AB $: $\cos\widehat B = \frac{BH}{AB}$. La relación depende sólo de
el ángulo (\cref{def:g8:cosine:def}), entonces
\[
\frac{AB}{a.\,C.} = \frac{BH}{AB} ;
\]
multiplicando ambos lados por $ a.\,C. \times AB $
(\cref{thm:g8:equations:balance}):
$ AB^2 = BH \times a.\,C. $.

\textbf{3.} los triangulos $ ABC $ (\omterm{def:g6:shapes:triangles}{en ángulo recto} en $ A $) y $ ACH $
(\omterm{def:g6:shapes:triangles}{en ángulo recto} en $ H $) comparte el ángulo $\widehat C $. Lado adyacente
sobre hipotenusa en cada uno:
\[
\cos\widehat C = \frac{AC}{CB} = \frac{CH}{AC}
\qquad\Longrightarrow\qquad
AC^2 = CH \times CB .
\]

\textbf{4.} $ AB^2 = BH \times a.\,C. $ lee $9 = BH \times 5$, entonces
$ BH = 1.8$; y $16 = CH \times 5$ da $ CH = 3.2$. Controlar:
$ BH + HC = 1.8 + 3.2 = 5 = a.\,C. $ --- como debe ser, ya que $ H $ se encuentra en
la hipotenusa entre $ B $ y $ C $.

\textbf{5.} Dividiendo cada relación por $ a.\,C. $:
\[
BH = \frac{AB^2}{a.\,C.},
\qquad
CH = \frac{AC^2}{a.\,C.} .
\]
Entonces $ BH $ y $ CH $ son \omterm{def:g6:propdata:def}{proporcional} a $ AB^2$ y $ AC^2$: el pie
de la altitud divide la hipotenusa en la proporción de los cuadrados
de las piernas ($1.8 : 3.2 = 9 : 16$ en la pregunta 4).

\textbf{6.} Sumando las dos relaciones y factorizando $ a.\,C. $
(distributividad):
\[
AB^2 + AC^2 = BH \times a.\,C. + CH \times a.\,C.
= (BH + CH) \times a.\,C. = a.\,C. \times a.\,C. = a.\,C.^2 .
\]
Esto es \cref{thm:g8:pythagoras:direct}, obtenido del coseno
solo --- una prueba genuinamente diferente de la de los cuatro triángulos \omterm{def:g6:measure:area}{área}
rompecabezas de \cref{exo:g8:pythagoras:11}.

\textbf{7.} Pitágoras en $ ABH $ (\omterm{def:g6:shapes:triangles}{en ángulo recto} en $ H $, hipotenusa
$[AB]$): $ AH^2 = AB^2 - BH^2$. Reemplazo $ AB^2$ por
$ BH \times a.\,C. $ (pregunta 2) y factorizando $ BH $:
\[
AH^2 = BH \times a.\,C. - BH^2
= BH \times (a.\,C. - BH)
= BH \times HC .
\]

\textbf{8.} Con las patas como base y altura, el \omterm{def:g6:measure:area}{área} de $ ABC $
es $\frac{AB \times AC}{2}$; con la hipotenusa $[a.\,C.]$ como base,
la altura es exactamente $ AH $, entonces el \omterm{def:g6:measure:area}{área} también es
$\frac{a.\,C. \times AH}{2}$ (\cref{thm:g7:areas:triangle}). equiparando
y duplicando: $ AB \times AC = a.\,C. \times AH $.

\textbf{9.} $ AH = \frac{AB \times AC}{a.\,C.} = \frac{3 \times 4}{5}
= 2.4$. Verificación de la pregunta 7:
$ AH^2 = 2.4^2 = 5.76$ y
$ BH \times HC = 1.8 \times 3.2 = 5.76$. Igual.

\textbf{10.} $ AH^2 = BH \times HC = 2 \times 8 = 16$, entonces
$ AH = \sqrt{16} = 4$ centímetro.

\textbf{11.} el punto $ A $ se encuentra en el círculo de \omterm{def:g4:geometry:circle}{diámetro} $[a.\,C.]$
(y difiere de $ B $ y $ C $), por lo que en el punto 2 de
\cref{thm:g8:cosine:circle} el triángulo $ ABC $ es \omterm{def:g6:shapes:triangles}{en ángulo recto} en
$ A $. Por construcción $(AH) \perp (a.\,C.)$ con $ H $ entre $ B $ y
$ C $: el \omterm{def:g6:lines:objects}{segmento} $[AH]$ es precisamente la altitud de la pregunta 7,
y $ AH^2 = BH \times HC = p \times q $.

\textbf{12.} $ AH^2 = 1 \times 2 = 2$: la figura construye, con
regla y compás solos, un \omterm{def:g6:lines:objects}{segmento} cuyo cuadrado es $2$ --- el
diagonal del cuadrado unitario, $\sqrt2 \approx 1.414$, de
\cref{ex:g8:pythagoras:diagonal}.

\textbf{13.} Receta: poner dos de punta a punta \omterm{def:g6:lines:objects}{segmentos} de longitudes
$ p = 1$ y $ q = n $; Dibuja el semicírculo cuyo \omterm{def:g4:geometry:circle}{diámetro} es el
entero \omterm{def:g6:lines:objects}{segmento} (longitud $ n + 1$); elevar el \omterm{def:g6:lines:perp}{perpendicular} en el
punto de unión; se encuentra con el semicírculo en un punto cuyo
la distancia al cruce es $\sqrt{1 \times n} = \sqrt n $. Para
$\sqrt5$: llevar $ p = 1$ y $ q = 5$ (un semicírculo de \omterm{def:g4:geometry:circle}{diámetro}
$6$).

\textbf{14.} $ A $ está en el círculo del centro $ O $, entonces
$ OA = \frac{p+q}{2}$, el \omterm{def:g3:shapes:shapes}{radio}. Si $ H \neq O $, el triángulo
$ AHO $ es \omterm{def:g6:shapes:triangles}{en ángulo recto} en $ H $ con hipotenusa $[OA]$, y una pierna es
más corto que la hipotenusa: $ AH < OA $. Si $ H = O $, entonces
$ AH = OA $ exactamente. en todos los casos $ AH \leq \frac{p+q}{2}$, con
igualdad precisamente cuando $ H $ es el centro --- es decir, cuando
$ p = q $.

\textbf{15.} cuadratura $ AH \leq \frac{p+q}{2}$ (ambos lados
positivo) y utilizando $ AH^2 = pq $:
\[
p \times q \leq \left(\frac{p+q}{2}\right)^2 ,
\]
con igualdad exactamente para $ p = q $. Reprueba algebraica, con la
identidades notables (\cref{pb:g8:equations:1}):
\[
\left(\frac{p+q}{2}\right)^2 - pq
= \frac{p^2 + 2pq + q^2 - 4pq}{4}
= \frac{p^2 - 2pq + q^2}{4}
= \left(\frac{p-q}{2}\right)^2 ,
\]
un cuadrado, por lo tanto nunca negativo --- y \omterm{def:g1:counting:zero}{cero} exactamente cuando
$ p = q$. El \omterm{pb:g8:cosine:1}{geométrico
        significar} de dos números nunca excede su
\omterm{def:g1:addition:def}{media suma}.
\end{solution}
